\documentclass[10pt,a4paper]{amsart}
\usepackage{amsmath,amssymb,amsthm,graphicx}
% Retained minimal preamble for the Stacks Project chapter more-morphisms.tex (More on Morphisms).
% Verbatim excerpt of the chapter preamble: Xy-pic setup (native preamble lines 14-18) and the native
% theorem declarations with the mathematical macros (native preamble lines 154-217). Copied mechanically;
% nothing added, removed or redefined. The wrapper class and the amsmath/amssymb/amsthm/graphicx base
% packages are supplied by the assembler, so the class/externaldocument part of the native preamble is omitted.
% For commutative diagrams we use Xy-pic
\usepackage[all]{xy}

% We use 2cell for 2-commutative diagrams.
\xyoption{2cell}
% Theorem environments.
%
\theoremstyle{plain}
\newtheorem{theorem}[subsection]{Theorem}
\newtheorem{proposition}[subsection]{Proposition}
\newtheorem{lemma}[subsection]{Lemma}

\theoremstyle{definition}
\newtheorem{definition}[subsection]{Definition}
\newtheorem{example}[subsection]{Example}
\newtheorem{exercise}[subsection]{Exercise}
\newtheorem{situation}[subsection]{Situation}

\theoremstyle{remark}
\newtheorem{remark}[subsection]{Remark}
\newtheorem{remarks}[subsection]{Remarks}

\numberwithin{equation}{subsection}

% Macros
%
\def\lim{\mathop{\mathrm{lim}}\nolimits}
\def\colim{\mathop{\mathrm{colim}}\nolimits}
\def\Spec{\mathop{\mathrm{Spec}}}
\def\Hom{\mathop{\mathrm{Hom}}\nolimits}
\def\Ext{\mathop{\mathrm{Ext}}\nolimits}
\def\SheafHom{\mathop{\mathcal{H}\!\mathit{om}}\nolimits}
\def\SheafExt{\mathop{\mathcal{E}\!\mathit{xt}}\nolimits}
\def\Sch{\mathit{Sch}}
\def\Mor{\mathop{\mathrm{Mor}}\nolimits}
\def\Ob{\mathop{\mathrm{Ob}}\nolimits}
\def\Sh{\mathop{\mathit{Sh}}\nolimits}
\def\NL{\mathop{N\!L}\nolimits}
\def\CH{\mathop{\mathrm{CH}}\nolimits}
\def\proetale{{pro\text{-}\acute{e}tale}}
\def\etale{{\acute{e}tale}}
\def\QCoh{\mathit{QCoh}}
\def\Ker{\mathop{\mathrm{Ker}}}
\def\Im{\mathop{\mathrm{Im}}}
\def\Coker{\mathop{\mathrm{Coker}}}
\def\Coim{\mathop{\mathrm{Coim}}}

% Boxtimes
%
\DeclareMathSymbol{\boxtimes}{\mathbin}{AMSa}{"02}

%
% Macros for moduli stacks/spaces
%
\def\QCohstack{\mathcal{QC}\!\mathit{oh}}
\def\Cohstack{\mathcal{C}\!\mathit{oh}}
\def\Spacesstack{\mathcal{S}\!\mathit{paces}}
\def\Quotfunctor{\mathrm{Quot}}
\def\Hilbfunctor{\mathrm{Hilb}}
\def\Curvesstack{\mathcal{C}\!\mathit{urves}}
\def\Polarizedstack{\mathcal{P}\!\mathit{olarized}}
\def\Complexesstack{\mathcal{C}\!\mathit{omplexes}}
% \Pic is the operator that assigns to X its picard group, usage \Pic(X)
% \Picardstack_{X/B} denotes the Picard stack of X over B
% \Picardfunctor_{X/B} denotes the Picard functor of X over B
\def\Pic{\mathop{\mathrm{Pic}}\nolimits}
\def\Picardstack{\mathcal{P}\!\mathit{ic}}
\def\Picardfunctor{\mathrm{Pic}}
\def\Deformationcategory{\mathcal{D}\!\mathit{ef}}
\title{The Mor functor from a finite free affine scheme into an affine scheme is representable by an affine scheme}
\author{The Stacks Project Authors}
\date{stacks-project:more-morphisms.tex:local-snapshot-20261011; GNU FDL 1.2 or later}
\begin{document}
\maketitle
\noindent\textit{Editorial note.} Statement span is the whole lemma environment lines 20099-20105; proof span is lines 20107-20168.\par\medskip
\noindent\textit{Editorial note.} The proof ends with the sentence 'Some details omitted.'; that closing sentence is retained verbatim. Under the project rule the main derivation must be complete and only routine verification may remain, which is the case here: the omitted part is the mechanical check that the assigned map P/J -> R' is well-defined and that the resulting T -> W has the required pullback property, and every ingredient for that check is already constructed in the retained text.\par\medskip
\noindent\textit{Editorial note.} The proof cites no label and no bibliography entry, so no external\_references are needed.\par\medskip
\noindent\textit{Editorial note.} The original writes the map name as b\_\{univ\} inside a math display and as b\_univ in surrounding prose; both are retained verbatim.\par\medskip
\noindent\textit{Editorial note.} extra\_packages is empty on purpose: the retained preamble already loads xy with options and a second declaration clashes.\par\medskip
\noindent\textit{Allowed standard prerequisites (editorial).}\begin{itemize}
\item Explicit prerequisite boundary. The functor Mor\_S(Z,X) of S-morphisms, finite free modules and bases, presentations Gamma(X,O\_X) = R[x\_i]/(f\_k), the Yoneda lemma, and pushforward/pullback of algebras along T -> S. The representing-object construction is retained in full.
\end{itemize}\medskip
\begin{lemma}
\label{lemma-hom-from-finite-free-into-affine}
Let $Z \to S$ and $X \to S$ be morphisms of affine schemes.
Assume $\Gamma(Z, \mathcal{O}_Z)$ is a finite free
$\Gamma(S, \mathcal{O}_S)$-module. Then $\mathit{Mor}_S(Z, X)$
is representable by an affine scheme over $S$.
\end{lemma}

\begin{proof}
Write $S = \Spec(R)$. Choose a basis $\{e_1, \ldots, e_m\}$
for $\Gamma(Z, \mathcal{O}_Z)$ over $R$. Choose a presentation
$$
\Gamma(X, \mathcal{O}_X) = R[\{x_i\}_{i \in I}]/(\{f_k\}_{k \in K}).
$$
We will denote $\overline{x}_i$ the image of $x_i$ in this quotient.
Write
$$
P = R[\{a_{ij}\}_{i \in I, 1 \leq j \leq m}].
$$
Consider the $R$-algebra map
$$
\Psi :
R[\{x_i\}_{i \in I}]
\longrightarrow
P \otimes_R \Gamma(Z, \mathcal{O}_Z), \quad
x_i \longmapsto \sum\nolimits_j a_{ij} \otimes e_j.
$$
Write $\Psi(f_k) = \sum c_{kj} \otimes e_j$ with $c_{kj} \in P$.
Finally, denote $J \subset P$ the ideal generated by the elements
$c_{kj}$, $k \in K$, $1 \leq j \leq m$. We claim that
$W = \Spec(P/J)$ represents the functor $\mathit{Mor}_S(Z, X)$.

\medskip\noindent
First, note that by construction $P/J$ is an $R$-algebra, hence
a morphism $W \to S$. Second, by construction the map
$\Psi$ factors through $\Gamma(X, \mathcal{O}_X)$, hence we obtain
an $P/J$-algebra homomorphism
$$
P/J \otimes_R \Gamma(X, \mathcal{O}_X)
\longrightarrow
P/J \otimes_R \Gamma(Z, \mathcal{O}_Z)
$$
which determines a morphism
$b_{univ} : W \times_S Z \to W \times_S X$.
By the Yoneda lemma $b_{univ}$ determines a
transformation of functors $W \to \mathit{Mor}_S(Z, X)$ which we
claim is an isomorphism. To show that it is an isomorphism it suffices
to show that it induces a bijection of sets
$W(T) \to \mathit{Mor}_S(Z, X)(T)$ over any affine
scheme $T$.

\medskip\noindent
Suppose $T = \Spec(R')$ is an affine scheme over $S$
and $b \in \mathit{Mor}_S(Z, X)(T)$. The structure morphism $T \to S$
defines an $R$-algebra structure on $R'$ and $b$ defines an $R'$-algebra map
$$
b^\sharp :
R' \otimes_R \Gamma(X, \mathcal{O}_X)
\longrightarrow
R' \otimes_R \Gamma(Z, \mathcal{O}_Z).
$$
In particular we can write
$b^\sharp(1 \otimes \overline{x}_i) = \sum \alpha_{ij} \otimes e_j$
for some $\alpha_{ij} \in R'$. This corresponds to an $R$-algebra map
$P \to R'$ determined by the rule $a_{ij} \mapsto \alpha_{ij}$. This
map factors through the quotient $P/J$ by the construction of the ideal
$J$ to give a map $P/J \to R'$. This in turn corresponds to a morphism
$T \to W$ such that $b$ is the pullback of $b_{univ}$.
Some details omitted.
\end{proof}

\end{document}
